% NB10c -- Arquitectura del modelo integrado: mecanismo 2D cuasi-estatico + termoquimica 0D en el tiempo
% + turbocompresor/intercooler cuasi-estacionarios. Flechas = informacion que se intercambia.
\documentclass[tikz,border=4mm]{standalone}
\usepackage{amsmath}
\usetikzlibrary{arrows.meta,calc,positioning}
\begin{document}
% Estilo comun de los diagramas TikZ del modulo VCR (NB10b) -- paleta del proyecto
\definecolor{manivela}{HTML}{2A78D6}
\definecolor{inferior}{HTML}{4A3AA7}
\definecolor{control}{HTML}{1BAF7A}
\definecolor{superior}{HTML}{EB6834}
\definecolor{tinta}{HTML}{3A3A38}
\definecolor{metal}{HTML}{C3C2B7}
\definecolor{perdida}{HTML}{E34948}
\definecolor{fuerza}{HTML}{C0392B}
\tikzset{
  barra/.style={line width=3.2pt, line cap=round},
  pin/.style={circle, fill=white, draw=tinta, line width=0.9pt, inner sep=1.6pt},
  fijo/.style={circle, fill=tinta, inner sep=2.6pt},
  vec/.style={-{Stealth[length=2.6mm]}, line width=0.9pt},
  F/.style={-{Stealth[length=3mm]}, line width=1.2pt, fuerza},
  cota/.style={{Stealth[length=1.8mm]}-{Stealth[length=1.8mm]}, line width=0.5pt, tinta},
  lbl/.style={font=\small, text=tinta},
}
\newcommand{\tierra}[1]{\fill[pattern=north east lines, pattern color=tinta] ($(#1)+(-5,-3)$) rectangle ($(#1)+(5,-6)$);
  \draw[tinta, line width=0.8pt] ($(#1)+(-5,-3)$) -- ($(#1)+(5,-3)$);}

\definecolor{mec}{HTML}{2A78D6}
\definecolor{ter}{HTML}{E34948}
\definecolor{qui}{HTML}{EB6834}
\definecolor{adm}{HTML}{1BAF7A}
\begin{tikzpicture}[x=1cm,y=1cm,
  bloque/.style={draw, line width=1.2pt, rounded corners=4pt, align=left, font=\small, inner sep=6pt, text width=6.6cm},
  info/.style={-{Stealth[length=3mm]}, line width=1.1pt, tinta},
  et/.style={font=\footnotesize, text=tinta, fill=white, inner sep=1.5pt, align=center}]
  \node[bloque, draw=adm, fill=adm!8] (A) at (0,4.6) {%
    \textbf{Admisi\'on} \hfill {\footnotesize cuasi-estacionario}\\
    compresor centr\'ifugo $\to$ intercooler $\varepsilon$-NTU\\
    $\Rightarrow p_3,\ T_3,\ \rho_3$, masa atrapada, $\lambda$};
  \node[bloque, draw=mec, fill=mec!8] (M) at (-4.6,0) {%
    \textbf{Mecanismo 2D} \hfill {\footnotesize cuasi-est\'atico}\\
    lazo vectorial + jacobiano: $y_D,\dot y_D,\ddot y_D$\\
    $V(\theta),\ \dot V(\theta)$ con $\omega$ = cte.\\[2pt]
    Newton-Euler: $F_O,F_C,F_D,N,T_O$};
  \node[bloque, draw=ter, fill=ter!8] (T) at (4.6,0) {%
    \textbf{Termoqu\'imica 0D} \hfill {\footnotesize EDO en $t$}\\
    zonas chorro/aire (Cantera ReactorNet)\\
    $\dot U=-p\dot V-\dot Q_w+\dot H_{f}$;\ \ $\dot Y_k=\dot\omega_k W_k/\rho$\\
    Woschni, inyecci\'on $\dot m_f(t)$, arrastre $\dot m_{arr}(t)$\\
    post: NO (Zeldovich ext.), holl\'in (Hiroyasu)};
  \node[bloque, draw=tinta, fill=metal!25] (E) at (0,-4.4) {%
    \textbf{Escape y desempe\~no}\\
    $T_5$, $\dot m_{ex}$ $\to$ turbina $\Rightarrow$ cierra $\dot W_T=\dot W_C$\\
    IMEP, BMEP (FMEP), par, potencia, BSFC, $C_s$};
  \draw[info] (M.north east) ++(-0.3,0) to[bend left=15] node[et, pos=0.55] {$V(t),\ \dot V(t)$\\ (tablas, $t=\theta/\omega$)} (T.north west);
  \draw[info] (T.south west) ++(0.3,0) to[bend left=15] node[et, pos=0.5] {$p(t)\Rightarrow F_{gas}$} (M.south east);
  \draw[info] (A.south west) to[bend right=10] node[et, pos=0.5] {$p_{adm}$} (M.north);
  \draw[info] (A.south east) to[bend left=10] node[et, pos=0.5] {$p_3,T_3,\lambda$\\ estado 1c} (T.north);
  \draw[info] (T.south) to[bend left=10] node[et] {$p,T$ en EVO} (E.north east);
  \draw[info] (M.south) to[bend right=10] node[et] {$T_O(\theta)$} (E.north west);
  \draw[info, dashed] (E.east) -- ++(5.6,0) |- node[et, pos=0.25] {$\Pi_C$\\(iteraci\'on\\turbo)} (A.east);
\end{tikzpicture}
\end{document}
